Cooperative multi-agent systems increasingly exchange learned messages, and
attention is a common way of deciding whose message matters
\citep{das2019tarmac,jiang2018atoc,jiang2018dgn}.\footnote{\LLMDisclosure}
Attention over $N$ agents is pairwise, so its cost grows quadratically, and
several methods restrict who attends to whom, citing communication or
computational cost among their reasons
\citep{jiang2018dgn,sun2020sparsecomm}. The computational
argument is an operation count: if each agent aggregates messages from $k$ peers
instead of $N-1$, the aggregation shrinks by a factor of about $N/k$.

This paper asks when that saving becomes elapsed time. Three costs have to be
separated, and an operation count conflates them. The first is \emph{what can
be removed at all}: a rule that selects peers by their attention scores needs
the scores, so every pairwise score has been computed by the time the selection
is made, and only the aggregation can shrink. The second is \emph{what the
implementation executes}: masking the discarded weights to zero leaves the
dense multiplication in place, while gathering the retained values removes it
and writes new intermediates instead. The third is \emph{what the hardware
spends its time on}: at the agent counts and head widths of this setting,
attention is memory-bound, so an implementation that removes arithmetic while
adding data movement can become slower. We study centralised inference on a
single device, the setting in which the methods we cite are trained and
evaluated, and not bandwidth-constrained distributed deployment, where reducing
transmitted messages is itself the objective.

\paragraph{Findings.} For \emph{score-then-select} rules, which form every score
before keeping the $k$ largest, at most half of the forward-pass attention
multiply--accumulates are removable at any sparsity, a limit noted before in a
memory-access model \citep{chen2023dfss}. Composed from library operators, such a rule is slower than dense attention at every agent count and sparsity we measure, the gathered form by up to thirty-seven times, because of memory traffic. A fused kernel
is faster than the best dense implementation only when one peer is retained (two
at a single configuration), at head widths $16$ and $32$, and without tensor
cores; in TF32 and bf16 the advantage is gone from $384$ agents upward. One
retained peer is a degenerate attention, since a softmax over one element is
identically one. Handing the fused kernel's own retained set to a kernel for free
shows where its time goes: from $384$ agents upward, at least nine tenths is
spent finding the set. Rules that \emph{select before scoring} escape the bound,
and attention over a pre-selected set is then up to sixteen times faster than
dense attention, but the selector decides the outcome: of an all-pairs, a grid
and a space-filling-curve selector, only the last reaches parity with dense
attention, at $512$ agents on best-case timings, and the ratio of medians
reverses that in one of our two sweeps.

\paragraph{Contributions.}
\begin{itemize}
  \item A cost model for peer selection (Section~\ref{sec:analysis}): a
    closed-form multiply--accumulate limit on score-then-select for the forward
    and backward passes (Proposition~\ref{prop:ceiling}); a selector-cost
    accounting that sorts selection mechanisms into four asymptotic classes
    (Proposition~\ref{prop:escape}, Corollary~\ref{cor:taxonomy}); and a byte
    criterion, $kd \le N$, that screens compaction before implementation.
  \item A measurement of what is executed (Section~\ref{sec:measurement}):
    masked, gathered and fused implementations of one function, verified
    equivalent, against the best of a vendor library, FlexAttention and a dense
    kernel searched over the sparse kernels' parameter space; three numerical
    formats; and a fixed-set comparison that separates the cost of finding the
    retained peers from the cost of using them.
  \item A characterisation of when sparse attention is faster in this regime:
    narrowly for score-then-select, and more broadly for selection before
    scoring when the selector is cheap in absolute terms and the remaining
    attention is fused (Sections~\ref{sec:fused}--\ref{sec:prescore}).
\end{itemize}

We release the kernels, the harness and the scripts that regenerate every table
and figure and check the quoted numbers. A pre-registered training study
(Appendix~\ref{app:training}) shows that the communicating architecture improves
return under the evaluated setup; it is not powered to rank sparsification
variants, and no claim in the main text rests on it. The measured envelope is
one GPU architecture, head widths $16$ to $512$ and at most $512$ agents
(Section~\ref{sec:limitations}).
